% =======================================================================================
% MODELO DE TCC - INSTITUTO FEDERAL DE SERGIPE (IFS)
% =======================================================================================
% Autores do Modelo (discentes de Ciência da Computação - IFS, Campus Itabaiana):
%   Mikael dos Santos Gois
%   Talyson Rodrigues Silva Nascimento
%   Kauan Gutierrez Santos Estácio
% Modelo Base: abnTeX2, adaptado para as normativas internas de TCC do IFS (Art. 31 e 34)
% Créditos dos modelos base (abnTeX2, DECOM/UFOP e Diogo D. Moreira) e das imagens: ver NOTICE
% Licença: CC BY-SA 4.0 (ver LICENSE)
% Este arquivo faz parte de um modelo opensource para estudantes do IFS.
% Não remova os créditos dos autores originais e do modelo base.
% Este modelo é fornecido "no estado em que se encontra", sem garantias de qualquer tipo.
% =======================================================================================

% [O que aparece no PDF Bookmark] {O que aparece impresso no documento}

\chapter[Resultados e Discussão]{RESULTADOS E DISCUSSÃO}
\label{sec:resultados-discussao}

% ATENÇÃO:
% Os números utilizados a seguir são fictícios.
% Substitua-os pelos resultados reais do TCC.

Este capítulo apresenta um exemplo de organização dos resultados obtidos em uma avaliação experimental, seguido de sua interpretação e discussão. Em um trabalho acadêmico, os resultados devem ser apresentados de forma objetiva, relacionados aos objetivos da pesquisa e analisados à luz da metodologia adotada. Os dados utilizados a seguir são fictícios e destinam-se exclusivamente à demonstração do uso do modelo.

\section{Apresentação dos Resultados}

Como exemplo, considera-se uma aplicação computacional avaliada em duas configurações:
uma sem otimização e outra com otimização.

Foram escolhidas duas métricas ilustrativas: tempo médio de execução, expresso em segundos,
e consumo médio de memória, expresso em megabytes.

A Tabela \ref{tab:resultados_exemplo}
apresenta os valores hipotéticos dessa comparação.

\begin{table}[htb]
    \centering
    \caption{Resultados ilustrativos de uma
    avaliação de desempenho}
    \label{tab:resultados_exemplo}

    \begin{tabular}{lrr}
        \hline
        \textbf{Configuração} &
        \textbf{Tempo médio (s)} &
        \textbf{Memória média (MB)} \\
        \hline
        Sem otimização & 4,8 & 620 \\
        Com otimização & 3,6 & 540 \\
        \hline
    \end{tabular}

    \fonte{Elaboração própria (dados fictícios
    para fins ilustrativos).}
\end{table}

No exemplo apresentado na Tabela
\ref{tab:resultados_exemplo}, o tempo médio
de execução passou de 4,8 para 3,6 segundos, correspondendo a uma redução de 25\%.

O consumo médio de memória passou de 620 para 540 MB, o que representa uma
redução aproximada de 12,9\%.

Essas porcentagens foram calculadas tomando os valores da configuração sem otimização
como referência.

\section{Discussão dos Resultados}

Os valores hipotéticos sugerem que a configuração com otimização teria melhor
desempenho nos dois critérios considerados, para as condições representadas neste exemplo.

Entretanto, uma comparação consistente exige descrever o ambiente de execução, as
características da carga de trabalho, a quantidade de repetições e a variabilidade
das medições.

Sem essas informações, não é possível avaliar a estabilidade das diferenças
observadas nem generalizar os resultados.

Na versão definitiva de um trabalho, esta discussão deve relacionar os achados aos
objetivos estabelecidos na Introdução, confrontá-los com as decisões metodológicas
e, quando pertinente, compará-los com estudos citados na Fundamentação Teórica ou nos
Trabalhos Correlatos.

\section{Limitações da Avaliação}

Os números utilizados neste exemplo não correspondem a experimentos reais.

Em uma pesquisa efetiva, devem ser explicitadas as limitações que afetam a
interpretação dos resultados, tais como o tamanho da amostra, as condições do ambiente
experimental, a representatividade dos cenários avaliados e os fatores que não
puderam ser controlados.

Dessa forma, as conclusões permanecem coerentes com as evidências obtidas e não
extrapolam o alcance da investigação.

% Exemplo de implementacao do indice de assuntos.
A avaliação experimental
\index{Avaliacao experimental@Avaliação experimental}
foi realizada para analisar o desempenho
\index{Desempenho} do sistema.

Foram consideradas métricas como tempo de
execução\index{Tempo de execucao@Tempo de execução}
e consumo de memória
\index{Consumo de memoria@Consumo de memória}.

% -------------------------------------------------------
% DICA PARA O ESTUDANTE — RESULTADOS E DISCUSSÃO
% -------------------------------------------------------
% Resultados e Discussão não devem ser apenas uma
% sequência de tabelas, gráficos e números.
%
% Primeiro, apresente objetivamente o que foi
% observado; depois, explique o que esses achados
% significam para o problema e os objetivos da
% pesquisa, relacionando-os à metodologia e à
% literatura utilizada.
%
% Ao adaptar o exemplo, o estudante deve substituir
% todos os dados fictícios pelas medições reais,
% identificar e explicar as métricas, comentar
% tendências ou diferenças observadas e, quando
% pertinente, comparar os achados com estudos
% anteriores.
%
% É importante discutir limitações e possíveis
% explicações sem tratar hipóteses como fatos nem
% afirmar resultados que os dados não comprovam.
